\documentclass[10pt,a4paper]{amsart}
\usepackage[margin=19mm]{geometry}
\usepackage{amsmath,amssymb,mathtools}
\usepackage[scr=rsfs,cal=euler]{mathalfa}
\usepackage{xfrac}
\usepackage[unicode,hidelinks,bookmarks=false]{hyperref}
\newcommand{\ie}{i.e.\ }
\newcommand{\cf}{cf.\ }
\newcommand{\resp}{respectively\ }
\newcommand{\Subsection}{\subsection}
\providecommand{\nobreakdash}{\nobreak\hbox{-}}
\setlength{\emergencystretch}{2em}
\multlinegap0pt
\usepackage{amssymb}
\usepackage{tikz}
\usepackage[normalem]{ulem}
\usepackage{braket}
\usepackage[shortlabels]{enumitem}
\newtheorem{prop}{Proposition}
\title{Associated varieties of integral minimal highest weight modules}
\author{Zhanqiang Bai, Jing Jiang, Rui Wang}
\date{Source: arXiv:2603.26371v2 (CC BY 4.0)}
\begin{document}
\maketitle

\noindent\emph{Source and licence.} Associated varieties of integral minimal highest weight modules. Version arXiv:2603.26371v2 (CC BY 4.0), distributed by arXiv under the Creative Commons Attribution 4.0 International licence, http://creativecommons.org/licenses/by/4.0/, as recorded in the arXiv licence metadata for this exact version and shown on the abstract page https://arxiv.org/abs/2603.26371v2. Full licence notice: \texttt{licenses/arxiv-2603.26371-CC-BY-4.0.txt}.\par\smallskip

\noindent\textit{Editorial scope.} Counted unit: the article's prop non-smooth (source0.tex lines 2225--2227). The article's complete original proof of it is delivered with it (source0.tex lines 2228--2273, one \texttt{proof} environment of 46 lines). No mathematics is written, repaired or re-derived: the statement and the proof are the article's own lines, in the article's own notation, and no retained line is substituted or re-flowed.  No further source-local context is needed: every cross-reference of every retained line resolves inside the two delivered spans.  Nothing was omitted from any retained span, so the package carries no omission record.  No retained line contains a citation, so the wrapper carries no bibliography and no bibliography entry.  No retained line contains a figure environment, an image-inclusion command or drawing code, so no image is delivered and the package declares no asset dependency.  Editorial formatting only, in addition: an AMS \texttt{amsart} preamble that restates the article's own declarations and macros used by the retained lines, namely the environments \texttt{prop} on separate counters, together with the standard packages those lines need.  The retained environments print in this document as Proposition 1 in order of appearance, whereas the article prints the same items with the numbering of its own full document.  The complete native source of the article as distributed in the arXiv e-print for this exact version is bundled unchanged as \texttt{sources/197273/source0.tex}.
\begin{prop}\label{non-smooth}
    For type $D_n$ $(n\geq 4)$, all elements of the following KL right cells are non-smooth: $\{w_0\mathcal{C}_i\mid 2\leq i \leq n-2\}$.
\end{prop}

\begin{proof}
{\bf Case 1.} Suppose that $n$ is even.

    Note that 
    \begin{align*}
        &w_0s_is_{i-1}\cdots s_k\\
        =&(-1,\cdots,-(n-i-1),-(n-i+1),\cdots,-(n-i+k+1),\\&-(n-i),-(n-i+k+2),\cdots,-n),
    \end{align*}
which has the pattern $\bar{1}\bar{3}\bar{2}$ for $1\leq k\leq i-1$.

For the permutation 
\begin{align*}
    &w_0s_is_{i+1}\cdots s_{n-2}s_n\\=&(n-i+1,1,-2,\cdots,-(n-i-1),-(n-i),\\&-(n-i+2),\cdots,-n),
\end{align*}
which has the pattern $31\bar{2}\bar{4}$.

 Note that 
    \begin{align*}
        &w_0s_is_{i+1}\cdots s_{i+(k-1)}\\
        =&(-1,\cdots,-(n-i-k),-(n-i+1),-(n-i-k+1),\\&-(n-i-k+2),\cdots,-(n-i),-(n-i+2),\cdots,-n),
    \end{align*}
which has the pattern $\bar{1}\bar{3}\bar{2}$ for $2\leq k\leq n-i$ and the pattern $\bar{2}\bar{1}\bar{3}$ for $k=1$.

{\bf Case 2.} Suppose that $n$ is odd.

  Note that 
    \begin{align*}
        &w_0s_is_{i-1}\cdots s_k\\
        =&(1,\cdots,-(n-i-1),-(n-i+1),\cdots,-(n-i+k+1),\\&-(n-i),-(n-i+k+2),\cdots,-n),
    \end{align*}
which has the pattern $1\bar{3}\bar{2}$ for $1\leq k\leq i-1$.   

For the permutation 
\begin{align*}
    w_0s_is_{i+1}\cdots s_{n-2}s_n=(&-(n-i+1),1,-2,\cdots,-(n-i-1),-(n-i),\\&-(n-i+2),\cdots,-n),
\end{align*}
 which has the pattern $\bar{2}\bar{1}\bar{3}$.

 Note that 
    \begin{align*}
        &w_0s_is_{i+1}\cdots s_{i+(k-1)}\\
        =&(1,\cdots,-(n-i-k),-(n-i+1),-(n-i-k+1),\\&-(n-i-k+2),\cdots,-(n-i),-(n-i+2),\cdots,-n),
    \end{align*}
which has the pattern $1\bar{3}\bar{2}$ for $2\leq k\leq n-i$ and the pattern $3\bar{1}\bar{2}\bar{4}$ for $k=1$.

\end{proof}

\end{document}
